% Articulation subsequent to generation of pi, phi, and e.
% Cumulative content; does not replace the preceding TRIT development.
\subsection{Exponential realizations of the generated coordinates}
\label{euler-trit-articulacion}

The regional coordinates already determined allow recognition of concrete
realizations of the TRIT quadratic collection. The order is substantive:
construction of the regions precedes evaluation of an exponential at
parameters \(\pi\) or \(2\log\varphi\). These expressions identify
the role of the coordinates in operators constructed from APP and TPK;
they do not participate in selecting the prefixes that produce them.

\subsubsection{Ternary cycle, nilpotent transport, and sheet incidence}

Multiplication \(a(r)=4r\) in \(\mathbb Z/9\mathbb Z\) has order
\(3\). On the units it determines the orbits \(\{1,4,7\}\) and
\(\{2,8,5\}\). In the augmentation module of either orbit, consisting
of integer combinations whose coefficients sum to zero, the basis
\((e_r-e_{4r},e_{4r}-e_{7r})\) gives
\[
 C=\begin{pmatrix}0&-1\\1&-1\end{pmatrix},
 \qquad C^2+C+I=0.
\]
The restricted form has Gram matrix
\(\bigl(\begin{smallmatrix}2&-1\\-1&2\end{smallmatrix}\bigr)\).
On residue evaluations modulo \(9\), the action distinguishes the sheets:
\(S(ax,ay)=aS(x,y)\), whereas \(P(ax,ay)=a^{-1}P(x,y)\).
The cyclic realization therefore preserves a contragredient orientation
between sum and product. Quotients of the integer lift are retained
in the corresponding arithmetic record.

Elementary parabolic transport is represented by
\[
 N=\begin{pmatrix}0&1\\0&0\end{pmatrix}.
\]
The areal--volumetric incidence already constructed in \eqref{eq:fav} gives
\[
 F_{\rm av}=\begin{pmatrix}0&1\\1&1\end{pmatrix},
 \qquad A=F_{\rm av}^2=\begin{pmatrix}1&1\\1&2\end{pmatrix}.
\]
The three operators act on declared real spaces: the augmentation module
tensored with \(\mathbb R\), the nilpotent-transport plane, and the
modal-incidence plane. Sharing a matrix dimension does not identify
these spaces or their coordinates.

\begin{proposition}[Concrete generators of the three types]
\label{euler-trit-generadores}
The operators
\[
 J_{\mathrm e}=\frac{2C+I}{\sqrt3},\qquad
 J_{\mathrm p}=N,\qquad
 J_{\mathrm h}=\frac{2A-3I}{\sqrt5}
\]
satisfy, respectively,
\[
 J_{\mathrm e}^{\,2}=-I,\qquad
 J_{\mathrm p}^{\,2}=0,\qquad
 J_{\mathrm h}^{\,2}=I.
\]
\end{proposition}
\begin{proof}
The relation for \(C\) implies \((2C+I)^2=-3I\).
Direct multiplication gives \(N^2=0\).
Finally, \(A\) has trace \(3\) and determinant \(1\), so
\(A^2-3A+I=0\) and \((2A-3I)^2=5I\).
\end{proof}

\begin{proposition}[Three evaluations of the exponential identity]
\label{euler-trit-evaluaciones}
In the preceding realizations,
\begin{equation}
 \boxed{
 \exp(\pi J_{\mathrm e})+I=0,\qquad
 \exp(J_{\mathrm p})=I+J_{\mathrm p},\qquad
 \exp(2\log\varphi\,J_{\mathrm h})=F_{\rm av}^{\,2}.}
 \label{euler-trit-tres-evaluaciones}
\end{equation}
\end{proposition}
\begin{proof}
The elliptic specialization of the tritic exponential gives
\(\cos\pi\,I+\sin\pi\,J_{\mathrm e}=-I\).
Nilpotency removes exactly the terms of degree at least \(2\)
in the second expression. For the third, \(\varphi^2=\varphi+1\) implies
\[
 \cosh(2\log\varphi)=\frac32,\qquad
 \sinh(2\log\varphi)=\frac{\sqrt5}{2}.
\]
Thus the exponential equals
\(\frac32I+\frac{\sqrt5}{2}J_{\mathrm h}=A\).
\end{proof}

The first equality represents the elliptic half-turn; the complete turn
satisfies \(\exp(2\pi J_{\mathrm e})=I\). The second retains a
nonzero linear term: the parabolic regime expresses transport, not absence
of evolution. The third realizes a unimodular hyperbolic advance. Its
eigenvalues are \(\varphi^2\) and \(\varphi^{-2}\); one direction
expands and the other contracts while preserving the determinant. The
three evaluations use different parameters of one common exponential collection.

In this collection, \(e\) enters through the scalar evaluation
\(\exp(I)=eI\) and the law \(\exp(sI)\exp(tI)=\exp((s+t)I)\).
Its role is not restricted to the parabolic type. The coordinate \(e\),
previously obtained from its region, is recognized here as the unit
evaluation of the exponential; this identity does not replace its
coinductive genealogy.

\subsubsection{Polynomial resolution of the regimes}

The three realizations admit a compact presentation on the direct sum
of their spaces. Define
\[
 \mathbb J=J_{\mathrm e}\oplus J_{\mathrm p}\oplus J_{\mathrm h},
 \qquad Q=\mathbb J^2.
\]
The preceding relations imply
\[
 \mathbb J^6=\mathbb J^2,\qquad Q^3=Q.
\]
Using \(Q\) keeps this operator square distinct from the
dodecaphase record \(K\).

\begin{proposition}[Polynomial regime projectors]
\label{euler-trit-proyectores}
The maps
\[
 p_{\mathrm e}=\frac{Q^2-Q}{2},\qquad
 p_{\mathrm p}=I-Q^2,\qquad
 p_{\mathrm h}=\frac{Q^2+Q}{2}
\]
are mutually orthogonal idempotents, sum to \(I\), and recover
the three summands of the representation.
\end{proposition}
\begin{proof}
The three polynomials take values \((1,0,0)\), \((0,1,0)\), and
\((0,0,1)\) at the points \(-1,0,+1\), respectively. The differences
between their squares and themselves, as well as their cross products,
are multiples of \(X^3-X\). Substituting \(Q\), which annihilates
that polynomial, proves the identities. On the three blocks, \(Q\)
acts as \(-I,0,I\).
\end{proof}

One thus obtains
\begin{align}
 \exp(t\mathbb J)
 ={}&p_{\mathrm e}(\cos t\,I+\sin t\,\mathbb J)
   +p_{\mathrm p}(I+t\mathbb J)\notag\\
   &+p_{\mathrm h}(\cosh t\,I+\sinh t\,\mathbb J).
 \label{euler-trit-exponencial-total}
\end{align}
Each summand preserves its quadratic relation, and conjugation reverses
its evolution. The total representation gathers the three regimes, but
does not by itself define a transition operator between them. Such a
transition requires the TPK transport maps with their domains, orientation,
and memory. The exponential articulation allows recognition of their
types without replacing those dynamics by an identification of fibers.

% FOCAL PROVENANCE (not printed)
% Recovered result: three generators and three exponential evaluations.
% Complete source owner:
% BASE/manuscrito/sections/hmt/09b_algebras_cuadraticas.tex:289-516.
% Causal antecedent of C and contragredient action:
% BASE/manuscrito/sections/hmt/03d_esqueleto_ciclotomico_app.tex:13-47,145-246.
% The three types and orientation are in:
% BASE/manuscrito/sections/hmt/02_trit_geometrias.tex:83-196.
% Recovered polynomial resolution:
% output/INVESTIGACION_HOLONOMIA_NONADICA_CRISTAL_APERIODICO_20260903/
% ALGEBRA_HOLONOMICA_UNIVERSAL_Y_EULER_TRITICO.md, sections 8-11 and 22.
% All four source owners were read fully. BASE denotes:
% /Users/ruben/Documents/HMT2/
% EL_CIERRE_HOLOGRAFICO_DEL_INFINITO_SUCESOR_102_CAPITULOS_2026-08-28_EN_TRABAJO/
% The identity P9(alpha)=delta from the specialized document is not incorporated:
% that identification does not belong to this block.
% Nor is a nilpotent direction identified with electronic vacancy,
% or a global transport representation by the direct sum presumed.
% CHECKS: standard Python, integers, and Fraction.
% C^2+C+I=0; (2C+I)^2=-3I; N^2=0; (2Fav^2-3I)^2=5I.
% 81 residue pairs verify the two contragredient actions.
% The 3 projectors were checked at Q=-1,0,+1; general proof in the text.
% 6 unique labels and balanced environments. No compilation.
